% !TeX root = ../papers/nc-p3.tex
\section{五次欠損の排除と階数十の二次曲線列}
\label{sec:rank-ten-quintic-centers}
\begin{theorem}\label{thm:rank-ten-quintic-centers}
反対階数10の曲面射影のない標準正則窓では、次数列 $(1,3,1)$ ばかりの分岐配置は存在しない。
従って分岐があれば全列の第二曲線は二次である。
\end{theorem}
\begin{proof}[二つの反対消去点]
この三次列の始終の中心を位置0の $p,z$ とする。もし $p=z$ なら最初の直線の異なる点から
同じ中心に戻る閉経路を反復できるので $p\ne z$。位置1,3の最初・最後の直線を $L,T$ とする。
$p$ の一意な前点 $p^-$ は $T$ の外にある。実際 $T$ の各点の後続は $z$ 一つである。
同様に $z$ の一意な後続 $z^+$ は $L$ の外にある。
位置2のねじれ三次上の順中心を $a$、逆中心を $b$ とする。前節より各々一つで、
両向きの中心が一致することはない。パラメータの対応する零点の一次式を $l_a,l_b$ とする。

$p$ の中央支持の消去点は $b$ であり、$z$ のそれは $a$ である。これを示す。
$p$ の正則三項鉛筆から、位置3の点を $p^-$ からの平面射影で読み、二つの行を縮約する。
$T$ の平面射影は直線であり、その一般点でこの二行は独立である。
実際正則三行三列鉛筆の可逆元を単位行列にすれば、従属方向はもう一元の左または右固有方向である。
固有空間が二次元なら非零な鉛筆元が階数1以下となってしまう。
これは定理\ref{thm:rank-nine-branch-exclusion}で用いた三項中心の連続純関係による排除で、
その階数下界は反対階数を使わない。従って従属方向は高々三点である。
二つの一次行の外積は $x_2(t)$ を $p$ の中央支持の消去点から射影した三座標に比例する。
この射の引戻し次数は2以下であり、中心を避ける三次曲線の射影ではない。
よって消去点はある $x_2(t)$ に等しい。逆向きの前節の三行行列はこの $t$ で
$p,z$ の両方に消され、階数2以下になる。逆向きでも末端直線への制限は同型なので、
ここが逆中心 $b$ である。順逆を交換して他方も従う。
\end{proof}
\begin{proof}[同じ五次式における階数一の矛盾]
$U=\ker p$ とし $R_0$ を $L$ に制限する。前節よりこの写像は $U\otimes H^0(L,\mathcal O(1))$ への同型である。
その逆写像の後で位置0を $p$ で評価した写像を $\beta$ とし、値域は $p$ の順支持 $S$（次元2）である。
$\beta_t(u)=\beta(u\otimes[-,t])$ と置くと
\[
             w(p,-,x_2(t),y)=\beta_t(A(t)y).
\]
$t_b$ では左辺は恒等零であり、一方 $A(t_b)$ は階数3で $U$ へ全射である。
ゆえに $\beta_t=l_b(t)\beta_0$、$\beta_0:U\to S$ は階数2の固定写像である。
$\ker\beta_0=ku$ と書くと、$p^-\to p$ から $A(t)p^-=h(t)u$ である。
$h$ は非零二次式である。恒等零なら一般の核の点 $x_3(t)$ が一定になる。
その根で $p^-$ が後続となるが、$p^-\notin T$ なので順中心を必要とする。
従って $h$ は $l_a^2$ の非零定数倍である。

$z(u)\ne0$ を確かめる。もし零なら $w(z,L,x_2(t),p^-)=0$ が全 $t$ で成立する。
$x_2$ は全位置2を張るので、$z$ の逆三項鉛筆を $p^-$ で評価した行列の位置1支持は
$L$ の消去二空間と $\ker z^+$ の交叉に入る。$z^+\notin L$ より交叉は一次元である。
しかし $p^-\notin T$ だからこの鉛筆元は非零で階数2以上であり矛盾する。

$T$ を張る二座標、$p^-$、さらに一補元を位置3の基底に取ると
$A(t)=(t_0L(t),t_1L(t),h(t)u,B(t))$ と書ける
（$L(t)$ はここだけ前節の一次列を表す）。五次欠損の式は
$h(t)\det(L(t),u,B(t))$ でその唯一の根は $t_a$、重複度5である。
よって二次因子 $h$ を除いても行列式は $t_a$ で消える。
この点では $w(z,a)=0$ だから $L(t_a),B(t_a)$ はどちらも $z$ に消される。
$z(u)\ne0$ より両列は従属であり、$h(t_a)=0$ なので $A(t_a)$ の階数が1以下になる。
これは順評価階数の下界2に反する。
\end{proof}

\begin{corollary}\label{cor:rank-ten-conic-center-cycles}
反対階数10で曲線分岐が存在するとすれば、各位相に順中心と逆中心が各一つあり、
それぞれが同じ向きの中心だけの四周期点列をなす。どの中心も特異三項または四次元で、
隣り合う二つの同向き中心がともに四次元になることはない。
各列の初めと終わりの三つの動く次数はともに $(1,2,1)$ であり、第二二次曲線は
その位相の順中心と逆中心をともに含む。
\end{corollary}
\begin{proof}
第二二次の最初の直線に乗る唯一の順中心へ写すと、その前点の一意性より中心間の単射ができる。
四位相を一巡すると有限集合の全単射なので、各位相に同数の中心があり、全てが次の中心として現れる。
逆向きも同じで、列の始終の対応で両向きの総数は等しい。
前節の二次曲線の鉛筆の議論より現れ得るのは特異三項または四次元で、四次元の直後は特異三項である。
そこで特異三項の順中心 $p$ を一つ選べる。その後続直線の外から来る次の点は
一つの平面 $b=0$ に入る。これは同鉛筆の末尾縮約に含まれる $S_p\otimes kb$ の純関係による。
$p$ と同位相に他の順中心があれば、その別の直線は $L_p$ と交わらず第二二次曲線はこの平面に入る。
さらにこの位置にある全ての列の最終直線はこの平面に入る。実際その前の曲線は
逆向きの第二二次曲線なので $L_p$ ではない。各位相に二中心以上あればこれら最終直線は二本以上あり、
先の別の二次曲線と合わせた平面内の総次数が4以上になって矛盾する。
従って各々一つである。二つの周期列は同じ位相で二方向中心が一致しないため互いに交わらない。

任意の順中心 $p$ を位置0、次の中心を $q$ とする。
第二二次の行列 $A(t)$ は $q$ の鉛筆である。$q$ が四次元なら $p$ は三次元で、
$D=A(t)V_3^*\subset U\otimes S_1$ は四次元、$\dim S_1D\le10-3=7$。
ゆえに $t_0u,t_1u$ の二列がある。一般階数2なら第一直線への制限像の消去式から
$A(t)$ の全列が $u' S_1+kv(t)$ に入る（前節の左一次関係の議論の一次行列版）ため
四列独立になれない。一般階数1も一次行列の全行または全列が固定方向となり不可能。
従って一般階数3で、第三曲線は二列の一次関係による直線である。

$q$ が特異三項なら $A(t)=u\alpha(t)^t+v(t)b^t$ の全非零値が階数2である。
第三曲線は $b=0$ 内に入り、同じ位置の最終直線も上の純関係によりこの平面に入る。
第三曲線が三次ならこの二成分の次数和に反する。二次なら $f_2=3-d_3=1$ で
第二曲線上の順中心 $r$ があり、その後続直線は階数2の $A(t_r)$ の核全体に等しいので
$b=0$ に入る。最初の後続直線と最終直線は異なる（各点の前点と次点がそれぞれ固定なら
動く位置が一つだけの列になる）。第三二次と二直線が同じ平面に入り再び矛盾する。
よって $d_3=1$ で、逆方向も同様である。
次数式により $f_2=2$、また $d_4=-2+b_2-f_3\ge0$（第三位置が最後なら $d_4=0$）。
従って第二曲線上には同じ位相の順中心と逆中心の双方がある。
\end{proof}

\begin{theorem}\label{thm:rank-ten-branch-exclusion}
反対階数10の標準正則窓で曲面射影がなければ、四辺の両射影の全ファイバーは零次元以下である。
従って同階数の全窓の正則性判定でなお未解決となり得るものは四辺が有限な窓に限られる。
\end{theorem}
\begin{proof}
前の系の分岐配置を仮定し、順側四周期列を $p_i$ とする。四次元の直後は特異三項なので、
巡回して $q=p_0$ を特異三項にできる。その中央支持の消去点を位置2の $d$ とする。
その鉛筆を先と同じく
\[
       A(t)=u(t_0\alpha_0+t_1\alpha_1)^t+(t_0v_0+t_1v_1)b^t
\]
とする。$\alpha_0,\alpha_1,b$ は位置2で $d$ を消す三文字、$u,v_0,v_1$ は位置3で
$q$ の前点 $p_3$ を消す三文字である。

一つ前の順中心 $p_3$ の列で、その第二二次曲線の上の順中心は $p_1$ である。
この二次曲線に沿う行列は、最初の直線上の順中心 $q$ の上記鉛筆であり全非零値が階数2である。
従って $p_1$ の後続直線はある $A(t)$ の核そのもので、$b=0$ の平面に入り $d$ を通る。
位置2における最終直線も $b=0$ に入る。実際その前の成分は二次曲線なので $q$ の後続直線ではなく、
$q$ の鉛筆が与える $S_q\otimes kb$ の純関係を使える。この二直線は異なる。

$q$ 自身の第二二次曲線はしたがって $b=0$ 内にはない。含まれれば別の二直線との次数和が4になる。
命題\ref{prop:singular-ternary-center-projections}を $q$ 自身の列に用いると、
第二二次曲線の $d$ からの射影は、第三直線の $p_3$ からの射影に線型同定される。
二次曲線の点からの射影は非定数であり、第三直線側の引戻し次数は高々1である。
よって $d$ は第二二次曲線上にある。
そのパラメータでの第一点は $p_1$ である。$d$ はこの中心の後続直線上にあり、前点は一意だからである。
一方第二二次曲線は $p_2$ を通り、その一意な前点も順側周期列により $p_1$ である。
第一直線のパラメータ表示は同型なので $d=p_2$ が従う。

ゆえに $w(p_0)$ の末尾縮約の三次元像は
$K=R_1\cap(V_1\otimes\ker p_2)$ に入り、その左支持は二次元の $S_q$ である。
$p_2$ の逆評価は階数3なので $\dim K=6-3=3$、この像は全 $K$ に等しい。
ところが $p_3$ の前点は $p_2$ なので、$w(p_3)$ の位置0での全縮約も $K$ に入る。
従ってその中央位置1での支持は $S_q$ に含まれ次元2以下となる。
これは系\ref{cor:no-binary-centers-without-two-lines}に反する。
こうして二つの分岐配置を全て排除した。
曲面以上なら定理\ref{thm:no-moving-plane-pencils}、正次元の曲線グラフなら
系\ref{cor:all-one-dimensional-point-graphs}が全次数判定を与えるので、有限辺だけが残る。
\end{proof}
